\section{Metodología}
\label{sec:metodologia}

\subsection{Objetivos}

\textbf{Objetivo general.} Diseñar e implementar un sistema que permita
registrar, priorizar, atender, cancelar y consultar solicitudes de atención
de forma eficiente, aplicando estructuras de datos apropiadas para cada
operación.

\textbf{Objetivos específicos.}
\begin{enumerate}
  \item Implementar una cola de prioridad formada por una cola FIFO por
    nivel de urgencia, cada una sobre una lista enlazada, para determinar
    la solicitud a atender en tiempo constante.
  \item Implementar la búsqueda y la cancelación de una solicitud pendiente
    por su identificador.
  \item Implementar una pila sobre lista enlazada que registre el historial
    de atenciones y permita, en una entrega posterior, deshacer la más
    reciente.
  \item Integrar las estructuras en una interfaz de consola que exponga las
    operaciones del sistema.
\end{enumerate}

\subsection{Alcance y limitaciones}

El sistema completo contempla: registrar una solicitud con un
identificador, un nombre y un nivel de urgencia (de 1, más urgente, a 5,
menos urgente); atender la solicitud de mayor prioridad; buscar y cancelar
una solicitud pendiente por identificador; consultar el historial de
atenciones; y deshacer la atención más reciente. Esta primera entrega
implementa todas estas operaciones excepto deshacer, que queda diseñada y
se implementará en la siguiente entrega.

Quedan explícitamente fuera del alcance del proyecto: persistencia en base
de datos, interfaz gráfica, acceso concurrente de múltiples usuarios y
autenticación. Se asume que los identificadores son proporcionados por
quien registra la solicitud; el sistema rechaza un identificador que ya
pertenezca a otra solicitud pendiente. Se asume también que la urgencia
toma solo cinco valores, supuesto del que depende el costo constante de la
cola de prioridad elegida.

\subsection{Requerimientos principales}

\begin{itemize}
  \item RF1: El sistema debe permitir registrar una solicitud con
    identificador único entre las pendientes, nombre y nivel de urgencia
    entre 1 y 5, rechazando valores fuera de ese rango.
  \item RF2: El sistema debe atender siempre la solicitud pendiente de
    mayor urgencia; en caso de empate, debe atenderse primero la que llegó
    antes.
  \item RF3: El sistema debe permitir buscar una solicitud pendiente por su
    identificador.
  \item RF4: El sistema debe permitir cancelar una solicitud pendiente de
    modo que nunca sea atendida.
  \item RF5: El sistema debe permitir deshacer la atención más reciente,
    devolviendo la solicitud correspondiente a la cola de pendientes.
  \item RF6: El sistema debe permitir consultar el historial de
    solicitudes atendidas.
\end{itemize}

\subsection{Diseño de la solución}

El sistema se organiza en los siguientes componentes: una clase
\texttt{Ticket} que representa una solicitud; una cola genérica
\texttt{LinkedQueue<T>} sobre lista enlazada; una cola de prioridad
\texttt{LevelPriorityQueue}, formada por un arreglo de cinco
\texttt{LinkedQueue<Ticket>}, una por nivel de urgencia; una pila genérica
\texttt{LinkedStack<T>}; y un historial \texttt{ServiceHistory} construido
sobre esa pila. Una clase fachada, \texttt{TriageSystem}, coordina la cola
de prioridad y el historial para que cada operación pública los mantenga
consistentes, y una clase \texttt{Main} ofrece el menú de consola.

\begin{figure}[h]
  \centering
  \fbox{\parbox{0.9\linewidth}{\centering
    \texttt{Main} (consola) \\[4pt]
    $\downarrow$ \\[4pt]
    \texttt{TriageSystem} \\[4pt]
    $\swarrow$ \hspace{3em} $\searrow$ \\[4pt]
    \texttt{LevelPriorityQueue} \hspace{1em} \texttt{ServiceHistory} \\[2pt]
    {\footnotesize 5 $\times$ \texttt{LinkedQueue<Ticket>}} \hspace{2.5em}
    {\footnotesize \texttt{LinkedStack<Ticket>}}
  }}
  \caption{Componentes principales del sistema y la estructura de datos
  que resuelve cada responsabilidad.}
  \label{fig:arquitectura}
\end{figure}

\subsection{Estructuras de datos seleccionadas}

\textbf{Cola de prioridad por niveles.} Representa el conjunto de
solicitudes pendientes. Está formada por un arreglo de cinco colas FIFO;
la cola de índice $u-1$ contiene, en orden de llegada, las solicitudes de
urgencia $u$. Debe soportar insertar una solicitud y extraer la más
prioritaria. Insertar consiste en encolar al final de la cola de su nivel,
y extraer consiste en recorrer los niveles desde el más urgente y desencolar
del primero que no esté vacío (Algoritmo~\ref{alg:atender}). Como el número
de niveles es la constante $k=5$, ambas operaciones cuestan $O(k)=O(1)$. El
desempate por orden de llegada (RF2) no requiere comparaciones adicionales:
lo garantiza el comportamiento FIFO de cada cola. Se descartó una lista
ordenada por prioridad, cuya inserción cuesta $O(n)$, y no se requiere un
\emph{heap}, cuyo costo $O(\log n)$ solo se justifica cuando la prioridad
puede tomar muchos valores distintos.

\textbf{Cola sobre lista enlazada con referencias a cabeza y cola.} Cada
nivel es una cola implementada como lista simplemente enlazada que mantiene
una referencia a la cabeza (frente) y otra al último nodo (final). La
referencia al final es la que permite encolar en $O(1)$; sin ella, encolar
requeriría recorrer la lista en $O(n)$. Se prefirió la lista enlazada sobre
un arreglo circular porque el número de solicitudes por nivel no se conoce
de antemano y porque la cancelación necesita retirar un elemento
arbitrario, operación que en una lista enlazada solo exige reenlazar el
nodo anterior con el siguiente, sin desplazar elementos.

\textbf{Búsqueda y cancelación por identificador.} La cola de prioridad
resuelve eficientemente ``\textquestiondown{}cuál es la siguiente solicitud a
atender?'', pero no ``\textquestiondown{}dónde está la solicitud con
identificador X?''. En esta entrega la búsqueda recorre las colas de los
cinco niveles, con costo $O(n)$ en el peor caso. La cancelación usa esa
búsqueda y luego elimina el nodo de la lista enlazada de su nivel en $O(1)$,
actualizando la referencia al final cuando el nodo eliminado es el último
(Algoritmo~\ref{alg:cancelar}); su costo total es $O(n)$. La misma búsqueda
permite rechazar identificadores repetidos al registrar, por lo que el
registro completo cuesta $O(n)$. Reducir estos costos a $O(1)$ esperado
mediante un índice por identificador es una de las mejoras previstas para
las siguientes entregas (Sección~\ref{sec:discusion}).

\textbf{Pila (lista enlazada simple, encabezado único).} Registra las
solicitudes ya atendidas en orden inverso de atención. Se eligió una pila
porque la operación prevista, deshacer la atención más reciente (RF5), es
exactamente el comportamiento LIFO: apilar al atender, en $O(1)$, y
desapilar al deshacer, en $O(1)$. Se implementó como lista enlazada con
solo referencia al tope, ya que todas las operaciones ocurren en ese
extremo. Consultar el historial completo (RF6) recorre la pila en $O(h)$,
siendo $h$ el número de atenciones.

En los algoritmos, $L[1..5]$ denota el arreglo de colas por nivel y $S$ la
pila del historial.

\begin{algorithm}[t]
\caption{Atender la siguiente solicitud}
\label{alg:atender}
\begin{algorithmic}[1]
\Function{ServeNext}{}
  \For{$u \gets 1$ \textbf{to} $5$}
    \If{$L[u]$ no está vacía}
      \State $t \gets$ \Call{Dequeue}{$L[u]$}
      \State apilar $t$ en $S$
      \State \Return $t$
    \EndIf
  \EndFor
  \State \Return \textsc{null}
\EndFunction
\end{algorithmic}
\end{algorithm}

\begin{algorithm}[t]
\caption{Cancelar una solicitud pendiente}
\label{alg:cancelar}
\begin{algorithmic}[1]
\Function{CancelRequest}{$id$}
  \For{$u \gets 1$ \textbf{to} $5$}
    \State $prev \gets$ \textsc{null}; $cur \gets L[u].head$
    \While{$cur \neq$ \textsc{null}}
      \If{$cur.value.id = id$}
        \If{$prev =$ \textsc{null}}
          \State $L[u].head \gets cur.next$
        \Else
          \State $prev.next \gets cur.next$
        \EndIf
        \If{$cur = L[u].tail$}
          \State $L[u].tail \gets prev$
        \EndIf
        \State \Return \textsc{true}
      \EndIf
      \State $prev \gets cur$; $cur \gets cur.next$
    \EndWhile
  \EndFor
  \State \Return \textsc{false}
\EndFunction
\end{algorithmic}
\end{algorithm}
